\section{Introduction}\label{sec:introduction}

A convex body is a compact convex subset of $\R^n$ with nonempty
interior. If $K=-K$, its center $0$ is an interior point, and its
polar body is
\[
 K^\circ=\{y\in\R^n:\langle x,y\rangle\le1
                         \text{ for every }x\in K\}.
\]
Write $|K|$ for $n$-dimensional Lebesgue volume and
$P(K)=|K|\,|K^\circ|$ for the volume product. The two determinant
factors cancel under an invertible linear change of coordinates, so
$P(TK)=P(K)$. Mahler's symmetric volume-product problem asks how small
this linear invariant can be. The conjectured value is $4^n/n!$,
the product of a cube and its polar cross-polytope. We prove this bound
in every dimension and classify the equality cases.

The extremal bodies are defined recursively. A centered nondegenerate
interval is a one-dimensional \emph{Hanner polytope}. Given Hanner
polytopes $H_1\subset\R^k$ and $H_2\subset\R^l$, their Cartesian
product and the convex hull
\[
 \begin{aligned}
 H_1\oplus_\infty H_2&=H_1\times H_2,\\
 H_1\oplus_1 H_2&=
 \conv\bigl((H_1\times\{0\})\cup(\{0\}\times H_2)\bigr)
 \end{aligned}
\]
are Hanner polytopes in $\R^{k+l}$. These are the $\ell_\infty$ and
$\ell_1$ sums, respectively. This class contains cubes, cross-polytopes,
and mixed recursive sums. We call an invertible linear image of a
Hanner polytope a \emph{linear Hanner body}.

\begin{theorem}\label{thm:main}
For every integer $n\ge1$ and every origin-symmetric convex body
$K\subset\R^n$,
\begin{equation}\label{eq:main-inequality}
 |K|\,|K^\circ|\ge\frac{4^n}{n!}.
\end{equation}
Equality holds if and only if $K$ is a linear Hanner body.
\end{theorem}

The linear equivalence in the statement respects the fixed origin
used for polarity.

\subsection{History and context}

Mahler introduced the volume-product problem in his work on convex
bodies and transference in the geometry of numbers~\cite{Mahler1939};
the planar inequality also goes back to him~\cite{MahlerPlanar1939}.
Reisner later characterized its equality cases as parallelograms,
as a consequence of his zonoid theorem~\cite[p.~340]{Reisner1986}.
Iriyeh and Shibata proved the complete three-dimensional symmetric
theorem, including equality~\cite[Theorem~1.1]{IriyehShibata}.
Fradelizi, Hubard, Meyer, Rold\'an-Pensado, and Zvavitch gave a
shorter proof based on equipartitions, together with another equality
argument and stability~\cite{FradeliziHubardMeyerRoldanZvavitch2022}.
Chen, Li, Xi, and Xu recently gave a shadow-flow proof of that theorem
in a preprint that also treats the nonsymmetric three-dimensional
problem~\cite[Theorem~7.1]{ChenLiXiXu2026}.

Hanner's recursively constructed spaces~\cite{Hanner1956} already
play a central role in the known higher-dimensional cases.
Saint-Raymond proved the sharp inequality for unconditional bodies
(those invariant under coordinate sign changes)~\cite{SaintRaymond1981};
Meyer and Reisner identified the Hanner equality cases in that
class~\cite{Meyer1986,Reisner1987}. Reisner also proved the sharp
inequality for zonoids, which are Hausdorff limits of Minkowski sums of
centered segments, and showed that equality within that class means a
parallelotope~\cite{Reisner1985,Reisner1986}; Gordon, Meyer, and
Reisner later gave a short geometric proof~\cite{GordonMeyerReisner1988}.
For local questions in Banach--Mazur distance, Nazarov, Petrov,
Ryabogin, and Zvavitch established strict local minimality, up to
linear equivalence, of the cube~\cite{NazarovPetrovRyaboginZvavitch2010}.
Kim extended this to all Hanner polytopes~\cite{Kim2014}.
Kim and Zvavitch obtained stability
in the unconditional class and the inequality in a neighborhood of
that class~\cite{KimZvavitch2015}.

A separate line of work seeks dimension-independent exponential
bounds without the sharp constant. Bourgain and Milman established
such a reverse Blaschke--Santal\'o inequality~\cite{BourgainMilman};
Kuperberg obtained an explicit stronger bound by a geometric argument
using Gauss linking integrals~\cite{Kuperberg}. Complex-analytic
approaches provide another setting for these questions. Nazarov used
H\"ormander's $\bar\partial$ theorem and Bergman kernels for the
reverse Blaschke--Santal\'o inequality~\cite{Nazarov2012}.
Berndtsson recast Kuperberg's proof using complex integrals
\cite{Berndtsson}. In a symplectic approach, Karasev proved the sharp
inequality for hyperplane sections and projections of $\ell_p$ balls
($1\le p\le\infty$) and Hanner polytopes~\cite[Theorems~4.2 and~4.4]{Karasev}.
These works give analytic and geometric precedents for volume-product
estimates; their estimates are not hypotheses of our proof.

The analytic ingredient is a conformal map of the disk for which
uniform angular measure on the boundary maps to a uniform imaginary
coordinate. It is a rotation of
Gross's example for the uniform distribution in his conformal Skorokhod
embedding~\cite[Section~3.2]{Gross}. A holomorphic mass estimate turns
that boundary law into a lower bound for an integral of feasible-simplex
volumes over the primal body. We prove the mass estimate directly by
Stokes' theorem, with its precise normalization; its general background
is the theory of generalized Lelong numbers~\cite[Sections~3 and~5]{Demailly}.
The symmetric-convex extremal-function framework of Lundin and Baran
was a further motivation for the analytic search; see
\cite{Lundin1985} and the formula recorded in
\cite[Proposition~1.2, pp.~246--247]{BosCalviLevenberg2001}.
The proof here instead uses the direct isolated-zero mass calculation,
not those extremal-function formulas as premises.
The equality argument leads to a classical property of normed spaces:
any three pairwise-intersecting closed balls have a common point.
Hanner studied the associated recursively constructed convex bodies
\cite{Hanner1956}. Lima's norm-additive decomposition criterion
\cite[Theorem~1.2]{Lima1978}, also recorded in
Hansen--Lima~\cite[Theorem~3.6(4)]{HansenLima1981}, is exactly the
metric-median condition used below. Reisner's unconditional equality
theorem already connects minimal volume product to this ball
intersection structure~\cite[Theorem~1]{Reisner1987}. Our endpoint
argument obtains the structure without assuming unconditionality;
Hansen--Lima's finite-dimensional classification is the final input
\cite[Corollary~7.4]{HansenLima1981}.

For a body without a prescribed center, the general Mahler problem
instead minimizes $|K|\,|(K-z)^\circ|$ over $z\in\operatorname{int}K$.
Its conjectured constant $(n+1)^{n+1}/(n!)^2$ is smaller than the
symmetric constant for $n\ge2$; a result for that general problem alone
therefore does not establish Theorem~\ref{thm:main}. The argument here
uses symmetry throughout and is independent of the general theorem
proved in the companion paper
\cite[Theorem~1.1]{OpenAIGeneralMahler2026}.

\subsection{Proof overview}

Section~\ref{sec:convex} records polarity and the two norm-sum formulas.
They show that every Hanner body has the stated volume product.
The remaining work is the lower bound and its converse equality
statement.

First let
\[
 A=\{X\in\R^d:|b_i\cdot X|\le1,\ 1\le i\le m\},
 \qquad A^\circ=\conv\{\pm b_1,\ldots,\pm b_m\},
\]
where the nonzero, pairwise nonproportional rows $b_i$ span $\R^d$.
Section~\ref{sec:lens} constructs
a planar lens with horizontal half-width $\lambda(t)$ at height $t$.
For $X\in\operatorname{int}A$, use this width to form
\[
 L_X=\{Y\in\R^d:|b_i\cdot Y|\le\lambda(b_i\cdot X),\ 1\le i\le m\}.
\]
Each independent $d$-tuple of signed constraints that are tight at
one point of $L_X$ determines a simplex: take the convex hull of $0$
and those signed rows. Let $\Sigma_X$ be the union of these simplices.
It lies in $A^\circ$.

Write $F$ for the conformal map from the unit disk to the lens.
For each set $I$ of $d$ independent rows, sample independent uniform angles
$\theta_i$ and solve $b_iZ=F(e^{i\theta_i})$, $i\in I$, for the
complex vector $Z$. Let $P_I$ be the probability that all remaining
row coordinates also lie in the closed lens. Section~\ref{sec:feasibility}
applies the mass estimate of Section~\ref{sec:holomorphic-mass} to high
powers of the inverse map: after a change of radial variables, the
Jacobian integrals concentrate on these boundary samples and force
$S:=\sum_I P_I\ge1$.
The uniform imaginary coordinate then identifies this probability sum
with the simplex integral in Section~\ref{sec:simplex}:
\[
 1\le S=\frac{d!}{4^d}\int_A|\Sigma_X|\,dX
                    \le\frac{d!}{4^d}|A|\,|A^\circ|.
\]
The first inequality is analytic; the last is the inclusion just
described. The exact identity behind this bound also controls the
missing volume $|A^\circ|-|\Sigma_X|$. Approximation extends the
inequality to all origin-symmetric convex bodies.

For equality, set $Q=K^\circ$ and pass to the body
\[
 B=K\oplus_1[-1,1],\qquad B^\circ=Q\times[-1,1]
\]
in dimension $d=n+1$. The norm-sum formula preserves equality at the
new dimension. Inner polytope approximations to $Q$ therefore have
integrated missing volume tending to zero. Section~\ref{sec:zero-deficit}
uses product neighborhoods to obtain, for each interior point of $B$
and each nonzero direction, a representation whose rows attain their
lens bounds at one vector satisfying all the constraints. This includes
degenerate limiting representations.
Finally, Section~\ref{sec:median} examines points approaching an apex
of $B$. The lens asymptotic gives a support-function inequality,
and convex separation gives a metric median for every triple in the
norm with unit ball $Q$: a point whose distances to each pair sum to
that pair's distance. The Hansen--Lima theorem then identifies $Q$,
and hence $K$, as a linear Hanner body
\cite[Corollary~7.4]{HansenLima1981}.

We next derive functional and entropy--transport consequences of the
all-dimensional theorem. The geometric proof begins in
Section~\ref{sec:convex} and ends in Section~\ref{sec:median}.

\subsection{Functional Mahler inequality}\label{sec:functional-mahler}

The volume-product problem has a functional form in which polarity is
replaced by convex conjugation. For a proper lower-semicontinuous convex
function $\varphi:\R^n\to\R\cup\{+\infty\}$, its Fenchel--Legendre transform is
\[
 \varphi^*(y)=\sup_{x\in\R^n}\{\langle x,y\rangle-\varphi(x)\},
 \qquad y\in\R^n.
\]
We use the convention $e^{-\infty}=0$.

\begin{corollary}[Even functional Mahler inequality]\label{cor:functional-mahler}
For every integer $n\ge1$ and every even proper lower-semicontinuous
convex function $\varphi:\R^n\to\R\cup\{+\infty\}$ such that
$0<\int_{\R^n}e^{-\varphi(x)}\,dx<\infty$,
\[
 \left(\int_{\R^n}e^{-\varphi(x)}\,dx\right)
 \left(\int_{\R^n}e^{-\varphi^*(y)}\,dy\right)\ge4^n.
\]
The second integral is allowed to be infinite, in which case the
inequality is immediate.
\end{corollary}

\begin{proof}
Set $f=e^{-\varphi}$, a nonzero integrable log-concave function. For
$z\in\R^n$, Fradelizi and Meyer use the polar function
\[
 f^z(y)=\inf_{\{x:f(x)>0\}}
        \frac{e^{-\langle x-z,y-z\rangle}}{f(x)}
\]
and the translation-minimized functional volume product
\[
 \mathcal P(f)=
 \left(\int_{\R^n}f(x)\,dx\right)
 \inf_{z\in\R^n}\int_{\R^n}f^z(y)\,dy.
\]
Their Proposition~1(a) states that the symmetric geometric Mahler
inequality in every dimension implies $\mathcal P(f)\ge4^n$ for every
even log-concave function on $\R^n$ with $0<\int_{\R^n}f<\infty$
\cite{FradeliziMeyer2008}. Theorem~\ref{thm:main} supplies this
all-dimensional hypothesis. Since $f^0=e^{-\varphi^*}$, the product
in Corollary~\ref{cor:functional-mahler} is at least $\mathcal P(f)$.
\end{proof}

The cited implication uses auxiliary convex bodies in dimensions
$n+m$ and lets $m$ tend to infinity; it is not an application only of
the geometric theorem in dimension $n$. The constant is sharp:
for $\varphi(x)=\sum_{i=1}^n|x_i|$, the transform $\varphi^*$ is zero on
$[-1,1]^n$ and $+\infty$ outside, and both integrals in
the corollary equal $2^n$. The equality classification
in Theorem~\ref{thm:main} concerns convex bodies. The limiting
implication does not by itself classify the even potentials for which
equality holds \cite[p.~1437]{FradeliziMeyer2008}.

For potentials without evenness, the companion paper on general convex
bodies gives the sharp lower bound $e^n$
\cite[Corollary~1.2]{OpenAIGeneralMahler2026}.

The functional inequality also has an entropy--transport formulation.
For a log-concave probability measure $\eta$ with a density, let
$H(\eta)=\int\log(d\eta/dx)\,d\eta$ denote its entropy relative to
Lebesgue measure, the negative of differential Shannon entropy.
For probability measures $\mu_1,\mu_2$ with finite first moments, set
\[
 \mathcal T(\mu_1,\mu_2)=
 \inf_{f\in\mathcal F}\left\{\int f\,d\mu_1+\int f^*\,d\mu_2\right\},
\]
where $\mathcal F$ is the class of proper lower-semicontinuous convex
functions $\R^n\to\R\cup\{+\infty\}$ and $f^*$ is the Fenchel--Legendre
transform. The cost $\mathcal T$ is allowed to be $+\infty$.

\begin{corollary}[Symmetric entropy--transport inequality]\label{cor:entropy-transport}
For every integer $n\ge1$, let $\eta_i(dx)=e^{-V_i(x)}\,dx$, $i=1,2$,
be full-dimensional log-concave probability measures, both symmetric
about the origin, with proper lower-semicontinuous convex potentials
$V_i$. Suppose their densities are essentially continuous, meaning that
$e^{-V_i}=0$ at $\mathcal H^{n-1}$-almost every point of
$\partial\operatorname{supp}\eta_i$. Their moment measures
$\nu_i=(\nabla V_i)_\#\eta_i$ have finite first moments,
$H(\eta_i)=-\int V_i\,d\eta_i$ is finite, and
\[
 H(\eta_1)+H(\eta_2)\le -n\log(4e^2)+\mathcal T(\nu_1,\nu_2).
\]
\end{corollary}

\begin{proof}
For $\rho_i=e^{-V_i}$, the identity
$|\nabla\rho_i|=|\nabla V_i|e^{-V_i}$ holds almost everywhere on the
interior of the support. Lemma~4 and the proof of Proposition~7 of
\cite{CorderoErausquinKlartag2015} give
$\int |\nabla\rho_i|<\infty$ and $V_i\in L^1(\eta_i)$, respectively.
The pushforward definition of $\nu_i$ therefore gives its finite first
moment, and $H(\eta_i)=-\int V_i\,d\eta_i$ is finite.
Gozlan's equivalence between the symmetric functional inverse Santal\'o
and entropy--transport inequalities, as stated in
\cite[Theorem~1.3]{FradeliziGozlanZugmeyer2021}, applies to
Corollary~\ref{cor:functional-mahler} with $c=4$.
\end{proof}
